\chapter{Discussion}
\label{ch:discussion}

The experimental results presented in Chapter~\ref{sec:results} demonstrate that the architectural placement and selection of preprocessing steps profoundly impact the viability of $z$-anonymity in smart metering. This chapter critically analyzes these findings, exploring the underlying causes of the observed behaviors and the trade-offs between privacy, data utility, and system performance.

\section{Privacy Challenges in High-Frequency Consumption Patterns}
The analysis of the input data distribution and the high suppression rates observed in the Baseline scenario highlight a fundamental conflict in high-frequency smart grid monitoring: the inverse relationship between peak energy demand and data availability. 

As shown in the analysis of data distribution (Section~\ref{sec:datadistribution}), there is a strong correlation between the average energy load and the entropy of the data stream. During the evening peak (16:00--22:00), the number of mathematically distinct values reaches its maximum (approximately 1,250 unique values per snapshot). This indicates that household behaviors are most distinct when energy consumption is highest. Because the $z$-anonymity algorithm operates by suppressing unique outliers, this high entropy triggers a massive wave of suppression exactly during the hours when the grid operator needs the data the most for load balancing and demand response. 

The near-total suppression observed in the Baseline at $z=50$ confirms that raw, high-precision data is inherently difficult to utilize. The exactness of the measurements scatters the data points so widely that finding $z$ identical values becomes statistically improbable. This proves that some form of data preprocessing is not just an option, but a strict requirement to achieve any functional level of data availability in a privacy-preserving grid.

\section{Evaluation of Information Loss Metrics for Skewed Data}
To address the inverse relationship between peak energy demand and data availability, Data Generalization (rounding) proved highly effective at increasing the Publication Ratio. However, the evaluation of the associated information loss revealed a critical flaw in how standard metrics interpret highly skewed IoT data.

While the standard Normalized Certainty Penalty ($\text{NCP}_{\text{std}}$) indicated an acceptable information loss of 10.8\% for integer rounding (Precision 0), this value is heavily diluted by the extreme outliers in the global data range (up to 9.26 kWh). When evaluated using the $\text{NCP}_{\text{eff}}$, which measures the interval width relative to the data span of the typical 95\% of users, the penalty skyrocketed to 118.5\%. 

This result exceeding 100\% is not a mathematical error, but rather a vital diagnostic indicator. It demonstrates that the resolution of the generalization mechanism (a 1.0 kWh interval) is physically coarser than the actual variance of the investigated population [0, 0.844 kWh]. Consequently, the distinguishability between individual households within this 95\% bracket is significantly reduced, as the vast majority of records are categorized into only two discrete states (0 or 1~kWh). While this transformation preserves a basic distinction between consumption levels falling below or above the 0.5~kWh rounding threshold, it results in a substantial loss of the original data variance. Therefore, we conclude that Precision 0 is too aggressive for this dataset, rendering the data analytically useless for typical consumption analysis. Precision 2 emerges as the optimal compromise, providing high availability ($>52\%$) while maintaining a functional effective utility limit.

\section{Impact of Temporal Aggregation on Bandwidth and Anonymity}
The Temporal Aggregation scenario yielded mixed results that challenge common assumptions about data smoothing. From a performance perspective, the strategy was highly successful, delivering deterministic bandwidth savings of up to 87.5\% for a 4-hour window. This makes it the premier choice for alleviating network congestion at the Edge.

However, contrary to the initial hypothesis that averaging would "smooth" out unique spikes and improve the Publication Ratio, the results showed a slight decrease in availability at higher $z$-thresholds compared to the Baseline. This occurred because calculating the arithmetic mean of high-precision floats generates new, synthetic floating-point values that are highly likely to be unique within a snapshot.

Despite failing to improve the quantitative Publication Ratio, Temporal Aggregation offers a significant, qualitative privacy benefit: \textit{irreversibility}. A 4-hour aggregated mean of 0.5 kWh can result from infinitely many combinations of eight individual 30-minute readings (e.g., a steady 0.5 kWh draw, or seven readings of 0.1 kWh and one massive spike of 3.3 kWh). Even if this aggregated value is published, an adversary learns far less about the exact, minute-by-minute behavioral routines of the household (such as exact meal times or sleep schedules) than they would from a raw baseline snapshot. Thus, aggregation provides robust behavioral privacy, even if it does not mathematically aid the $z$-anonymity clustering process.

\section{Impact of Local Suppression on Global Data Availability}
The Local Prefiltering scenario evaluated the effectiveness of shifting the privacy check to the Fog Layer. As described in Section~\ref{sec:scenarios}, this distributed architecture relies on an initial generalization step at the Edge Layer ($p=2$). This preparation was a strict prerequisite; without it, the high resolution of the raw data ($p=3$) combined with the small neighborhood size (100 users per gateway) led to extremely high suppression rates during preliminary testing.

However, even with the data generalized to two decimal places, the results show a severe degradation of the global Publication Ratio, despite excellent bandwidth savings (up to 99\% at Local $z=10$). This severe drop is caused by the restricted view of the individual gateways. Because each Gateway only monitors its own local neighborhood, it lacks knowledge of the global population. A specific energy reading might be common across the entire city of 5,529 users, but appear uniquely rare within a single street. Consequently, the Gateway aggressively suppresses these local outliers before they ever reach the Central Entity.


This behavior highlights a fundamental trade-off in distributed IoT privacy: balancing message suppression against reduced data precision. To make Fog-based prefiltering viable, the architecture must rely on lowering numerical precision at the data source in order to reduce sparsity at the neighborhood level. The results suggest that, for distributed $z$-anonymity to function effectively, transmitting generalized data that satisfies local thresholds is preferable to high-precision reporting that leads to excessive data loss. Without this compromise, the limited local perspective of the Gateway becomes a bottleneck that significantly degrades global data utility.

\section{Limitations of the Evaluation}

While this simulation provides clear architectural insights, certain limitations regarding the experimental design must be acknowledged. First, the evaluation relies on the idealized assumption that all 5,529 smart meter readings for a given interval arrive at the Central Entity simultaneously. In a real-world deployment, data transmission is rarely instantaneous, measurements often arrive slightly earlier or later than intended. Because the Snapshot $z$-anonymity model used in this research enforces rigid 30-minute processing windows, tuples that arrive late would effectively be shifted into the subsequent snapshot. This temporal mismatch could alter the experimental results. A late-arriving tuple is forced to be evaluated against a different "crowd" from a different time period. Since energy consumption patterns are highly time-dependent, a value that was common in its original timeframe might be an outlier in the next, potentially leading to higher suppression rates than those observed in this perfectly synchronized simulation. While such increased suppression would reduce the overall data utility, it simultaneously demonstrates the effectiveness of the $z$-anonymity model. The algorithm correctly identifies and blocks tuples that have become unique due to temporal shifts, thereby maintaining the privacy guarantee even under non-ideal network conditions.

Second, the Fog Layer simulation logically partitioned users into static, arbitrary neighborhoods of 100 households based on their assigned identifiers. In a real-world deployment, physical gateways might aggregate households with shared demographic or geographic characteristics. A random assignment fails to capture natural local clusters of similar behavior, which likely contributed to the aggressively high local suppression rates observed in Section~\ref{sec:localprefiltering}.

Finally, the evaluation relies on the London Households dataset recorded between 2011 and 2014 (with a subset from 2012 used in this study). Modern smart grids are increasingly characterized by the proliferation of low-carbon technologies~\cite{DataVariability}. This includes both massive new consumption loads, such as Electric Vehicles (EVs) and heat pumps, as well as localized energy generation from residential solar panels~\cite{NewConsumptionLoads}. These technologies introduce new, highly distinct peaks and troughs in consumption patterns. Consequently, the baseline sparsity and entropy observed in this thesis may actually underestimate the privacy challenges present in contemporary power grids, where advanced load monitoring algorithms can accurately profile individual appliance usage and user behavior~\cite{NILM}.


\section{Future Work}
While this research provides a structural framework for distributed privacy preprocessing, several avenues for future work remain. 

First, future studies should evaluate the impact of varying the temporal window size~($\Delta t$). Investigating how different snapshot intervals (e.g., 1-hour or 2-hour boundaries) influence the "boundary effect" and the resulting data sparsity would help optimize the baseline configuration for $z$-anonymity. 


Second, while this thesis focused on deterministic preprocessing techniques like scalar rounding and temporal aggregation, future research could compare these results against probabilistic frameworks. Specifically, despite the theoretical challenges regarding noise injection and privacy budget depletion discussed in Section~\ref{sec:FormalPrivacyModels}, a comparative study using Differential Privacy at the edge could be conducted to quantify the actual utility impact of adding random noise versus reducing data precision. Such a study could determine if probabilistic approaches offer a different utility-privacy balance for high-frequency energy data.

Finally, the performance of the Fog Layer could be optimized through behavior-aware network topologies. Instead of arbitrary partitioning, future architectures could dynamically cluster Edge nodes into Fog gateways based on historical usage similarity. Investigating whether grouping households with similar routines natively increases the success rate of local prefiltering—without requiring complex inter-gateway communication—presents a highly promising direction for decentralized IoT privacy.